% Mục 4.1. Nguồn: include/lock.h, backoff.hpp, platform.hpp. Lệnh AArch64: c++ -O2 -S (Apple clang 21).
Mọi lock trong \texttt{include/lock.h} có đúng \texttt{acquire()} và \texttt{release()} không tham số như
khung của đề. Với khai báo \texttt{using pc::lock\_mcs;}, dòng \texttt{lock\_mcs lock;} của khung đề
biên dịch nguyên văn. Lấy lock dùng acquire, nhả lock dùng release. Cặp này làm mọi phép ghi trong critical section trước nhìn thấy
được ở chủ lock sau. Trạng thái lock được khai báo \texttt{alignas(kCacheLine)}, nên mỗi lock chiếm riêng
một cache line (128 byte trên hai máy Apple, 64 byte trên \PCname{LAP1184}). Nhờ vậy không có false
sharing, tức hai biến độc lập nằm chung một line và làm mất hiệu lực line của nhau.
% nguồn: include/platform.hpp, kCacheLine = 128 khi __APPLE__ && __aarch64__, 64 ở nơi khác.

TAS lặp \texttt{exchange(true, acquire)} (lệnh \texttt{swpab}) tới khi đọc được \texttt{false}. TAS nhả lock bằng
\texttt{store(false, release)} (\texttt{stlrb}). TTAS chờ bằng \texttt{load(relaxed)} (\texttt{ldrb}) rồi
mới \texttt{exchange}. Phép đọc chỉ là gợi ý, \texttt{exchange} thắng mới là phép acquire.

Mọi vòng spin gọi \texttt{cpu\_relax()} giữa hai lần thử, tức là \texttt{pause} trên x86 và \texttt{isb}
trên AArch64. Lệnh \texttt{isb} được chọn vì \texttt{yield} của Apple gần như không làm gì (0,3~ns so với
12,1~ns của \texttt{isb}, cùng phép đo). \texttt{isb} dừng pipeline trong khoảng thời gian đó, nên
thread chờ đọc lock thưa hơn.
% nguồn: tables.tex, bảng sleeplatency: yield_insn trung vị 0,300 ns; isb trung vị 12,1 ns.

MCS~\cite{mellorcrummey1991mcs} giữ \texttt{tail\_} (\texttt{nullptr} là rảnh). Mỗi thread chờ trên cờ
\texttt{locked} trong nút của mình, nên mỗi lần chuyển giao chỉ ghi một line của một thread, theo thứ tự
FIFO. Trong các đoạn mã dưới đây, \texttt{relaxed}, \texttt{acquire}, \texttt{release},
\texttt{acq\_rel} viết tắt cho \texttt{std::memory\_order\_*}.
\begin{lstlisting}
me->next.store(nullptr, relaxed); me->locked.store(true, relaxed);  // not yet reachable
pred = tail_.exchange(me, acq_rel);                 // release: publish node; acquire: take lock
if (pred == nullptr) return;
pred->next.store(me, release);
while (me->locked.load(acquire)) cpu_relax();       // spin on own cache line
// release(): succ = next.load(acquire); if null, CAS tail_ me -> null (release);
//            on failure wait until next != null (acquire)
succ->locked.store(false, release);                 // hand over
\end{lstlisting}
Hai phép ghi relaxed an toàn vì nút chỉ trở nên với tới được qua hai phép release theo sau. Nếu CAS
\texttt{tail\_} thất bại, một thread đã đổi \texttt{tail\_} nhưng chưa nối vào \texttt{next}. Chủ lock
phải chờ nó, vì nếu bỏ đi thì thread đó chờ mãi. Vòng chờ này ngắn, vì thread kia đã \texttt{exchange}
xong và lệnh kế tiếp của nó là ghi \texttt{next}. Một thread chỉ vào critical section khi nó thấy
\texttt{tail\_} rỗng hoặc khi thread trước nó đặt \texttt{locked} của nó về \texttt{false}, nên tại mỗi
lúc chỉ có một chủ lock. Hàng đợi theo thứ tự FIFO, nên không thread nào bị
đói. Trên AArch64 các thao tác thành \texttt{swpal}
(\texttt{exchange} lên \texttt{tail\_}), \texttt{stlr} (nối \texttt{next}), \texttt{ldaprb} (chờ
\texttt{locked}), \texttt{ldapr} (\texttt{release()} đọc \texttt{next}), \texttt{casl} (CAS
\texttt{tail\_}) và \texttt{stlrb} (trao lock). Nút \PCname{mcs_node} cũng chiếm riêng một line.
% nguồn: c++ -O2 -S của mcs_acq/mcs_rel (Apple clang 21): swpal, stlr, ldaprb; ldapr, casl, stlrb.

Giới hạn thứ nhất của MCS là một thread không giữ được hai lock MCS cùng lúc, vì nút là biến
\texttt{thread\_local} dùng chung cho mọi \PCname{lock_mcs} do giao diện không có tham số. Giới hạn thứ
hai là lock được giao cho đúng một thread, nên nếu thread đó bị preempt thì mọi thread phía sau phải chờ
scheduler.

\PCname{lock_tas_bo<B>} và \PCname{lock_ttas_bo<B>} gọi chính sách \texttt{B} sau mỗi lần thử thất bại.
TTAS chỉ lùi lại khi thua \texttt{exchange} (BackoffLock~\cite{herlihy2020art}). Memory order giống TAS và
TTAS, nên lập luận về tính đúng ở trên giữ nguyên. \texttt{B} là một trong
ba chính sách. \PCname{sleep} theo mẫu của đề, gọi \texttt{sleep\_for} với \texttt{delay} nhân đôi từ
100~ns tới \texttt{MAX\_DELAY} $= 100\,000$~ns. \PCname{spin} quay ngẫu nhiên, với cửa sổ nhân đôi từ 4
tới 1024 lần \texttt{cpu\_relax}. \PCname{yield} quay 64 lần rồi gọi \texttt{yield}. \texttt{sleep\_for}
là system call. Trên \PCname{apple-m5}, yêu cầu 100~ns ngủ trung vị 2\,670~ns (nhóm đo bằng
\texttt{bench/sleep\_latency.cpp}).
% nguồn: tables.tex, bảng sleeplatency, sleep_for 100 ns: trung vị 2670 ns.

Lock thư viện được bọc inline sau cùng giao diện. Đó là \texttt{std::mutex}, \texttt{pthread\_mutex\_t},
\texttt{os\_unfair\_lock} (macOS), \texttt{pthread\_spinlock\_t} (Linux), \texttt{omp\_lock\_t} (khi có
OpenMP). Bảng~\ref{tab:impl:names} (phụ lục) liệt kê tên mọi cài đặt và máy có chúng.
